\section{Les boucles imbirquées}

\begin{UPSTIexercice}{Trace d'une boucle imbriquée}
    On exécute un programme contenant les lignes suivantes : 
    \begin{lstlisting}[language=C]
for (int ligne = 1; ligne <= 2; ligne = ligne + 1)
{
    for (int colonne = 1; colonne <= 3; colonne = colonne + 1)
    {
        printf("(%d,%d) ", ligne, colonne);
    }
    printf("\n");
}
\end{lstlisting}
    \UPSTIquestion{Compléter le tableau de trace ci-dessous, en indiquant, pour chaque passage dans le bloc le plus interne, la valeur de \texttt{ligne}, celle de \texttt{colonne}, et ce qui est affiché.}

    \begin{center}
        \begin{tabular}{|c|c|c|c|}
            \hline
            \textbf{Passage n°} & \texttt{ligne} & \texttt{colonne} & \textbf{Affiché} \\
            \hline
            1 & & & \\[0.3cm]
            \hline
            2 & & & \\[0.3cm]
            \hline
            3 & & & \\[0.3cm]
            \hline
            4 & & & \\[0.3cm]
            \hline
            5 & & & \\[0.3cm]
            \hline
            6 & & & \\[0.3cm]
            \hline
        \end{tabular}
    \end{center}
    \UPSTIquestion{Combien de fois le \texttt{printf("\textbackslash n")} (celui qui n'est \textbf{pas} imbriqué dans la boucle \texttt{colonne}) est-il exécuté au total ? Justifier.}
    \UPSTIquestion{Donner alors l'affichage complet du programme dans la console.}
\end{UPSTIexercice}

\begin{UPSTIprofOnlyEnv}
    \begin{UPSTIcorrectionP}{Trace d'une boucle imbriquée}
        \begin{center}
            \begin{tabular}{|c|c|c|c|}
                \hline
                \textbf{Passage n°} & \texttt{ligne} & \texttt{colonne} & \textbf{Affiché} \\
                \hline
                1 & 1 & 1 & \texttt{(1,1) } \\
                \hline
                2 & 1 & 2 & \texttt{(1,2) } \\
                \hline
                3 & 1 & 3 & \texttt{(1,3) } \\
                \hline
                4 & 2 & 1 & \texttt{(2,1) } \\
                \hline
                5 & 2 & 2 & \texttt{(2,2) } \\
                \hline
                6 & 2 & 3 & \texttt{(2,3) } \\
                \hline
            \end{tabular}
        \end{center}
        Le \texttt{printf("\textbackslash n")} externe s'exécute \textbf{2 fois} : une fois par tour de la boucle \texttt{ligne}, une fois que la boucle interne \texttt{colonne} est entièrement terminée pour cette ligne.
        L'affichage complet du programme est donc :
        \begin{lstlisting}[language=bash,style=console]
(1,1) (1,2) (1,3)
(2,1) (2,2) (2,3)
\end{lstlisting}
    \end{UPSTIcorrectionP}
\end{UPSTIprofOnlyEnv}


\section{\texttt{do...while} : au moins une fois}

\begin{UPSTIinfor}{Rappel : la particularité du \texttt{do...while}}
    Une boucle \texttt{do...while} teste sa condition \textbf{après} avoir exécuté le bloc : celui-ci s'exécute donc toujours \textbf{au moins une fois}, même si la condition est fausse dès le départ. C'est la seule différence de fond avec un \texttt{while}, où la condition est testée \textbf{avant} toute exécution du bloc.
\end{UPSTIinfor}

\begin{UPSTIexercice}{\texttt{while} contre \texttt{do...while} : quand ça se voit}
    On considère les deux programmes suivants, tous deux exécutés avec \texttt{stock} valant initialement \texttt{0}.
    \noindent\\
    \begin{minipage}[t]{0.48\linewidth}
        \textbf{Programme A}
        \begin{lstlisting}[language=C]
int stock = 0;
while (stock > 0)
{
    printf("Stock : %d\n", stock);
    stock = stock - 1;
}
printf("Fin A\n");
        \end{lstlisting}
    \end{minipage}\hfill
    \begin{minipage}[t]{0.48\linewidth}
        \textbf{Programme B}
        \begin{lstlisting}[language=C]
int stock = 0;
do
{
    printf("Stock : %d\n", stock);
    stock = stock - 1;
} while (stock > 0);
printf("Fin B\n");
        \end{lstlisting}
    \end{minipage}\\
    \UPSTIquestion{Prédire tout ce qu'affiche le programme A.}
    \UPSTIquestion{Prédire tout ce qu'affiche le programme B.}
    \UPSTIquestion{Expliquer les éventuelles différences.}
\end{UPSTIexercice}

\begin{UPSTIprofOnlyEnv}
    \begin{UPSTIcorrectionP}{\texttt{while} contre \texttt{do...while} : quand ça se voit}
        Programme A : la condition \texttt{stock > 0} est testée \textbf{avant} tout passage dans le bloc ; avec \texttt{stock = 0}, elle est fausse immédiatement, donc le bloc ne s'exécute jamais. Affichage : \texttt{Fin A}.

        Programme B : le bloc s'exécute une première fois \textbf{avant} tout test. Affichage :
        \begin{lstlisting}[language=bash,style=console]
Stock : 0
Fin B
        \end{lstlisting}
        (après ce premier passage, \texttt{stock} devient \texttt{-1}, et \texttt{-1 > 0} est faux, donc la boucle s'arrête.)

        Les deux programmes n'affichent \textbf{pas} la même chose : avec une condition initialement fausse, le \texttt{do...while} exécute quand même une fois son bloc, contrairement au \texttt{while}. C'est exactement le cas où la différence entre les deux structures devient visible.
    \end{UPSTIcorrectionP}
\end{UPSTIprofOnlyEnv}
